\section{Introduction}

Parton distribution functions (PDFs) are key quantities for gaining an understanding of hadron structure and for making predictions for the cross sections in high-energy scattering experiments. In QCD factorization theorems for hard scattering processes~\cite{Collins:1989gx}, the relevant PDFs are defined in terms of the nucleon matrix elements of light-cone correlation operators. For example, in dimensional regularization with $d=4-2\epsilon$, the bare unpolarized quark PDF is
\begin{align} \label{eq:lcpdf}
	q(x,\epsilon) \equiv\! \int\! {d\xi^-\over 4\pi} e^{-ixP^+\xi^-}\! \langle P | \bar{\psi}  (\xi^-) \gamma^+ U (\xi^-, 0) \psi (0) | P \rangle ,
\end{align}
where $x$ is the momentum fraction, the nucleon momentum $P^\mu=(P^0,0,0,P^z)$, $\xi^{\pm}=(t\pm z)/\sqrt{2}$ are the light-cone coordinates, and the Wilson line is
\begin{align}
	U(\xi^-,0) = P\exp\biggl(-ig\int_0^{\xi^-} d\eta^- A^+(\eta^-)\biggr)\ .
\end{align}
Most often the bare PDF is renormalized in the $\overline{\ensuremath{\operatorname{MS}}}$ scheme to obtain $q(x,\mu)$, and this renormalized PDF  is used to make predictions for experiment. The relation is
\beq
	q(x,\epsilon) = \int_x^1 {dy\over y}\:  Z^{\overline{\ensuremath{\operatorname{MS}}}}\Bigl({x\over y},\epsilon,\mu\Bigr) q(y,\mu)\ ,
\eeq
where $\mu$ is the renormalization scale, and we have suppressed the flavor indices in the renormalization constant $Z^{\overline{\ensuremath{\operatorname{MS}}}}$ and PDFs. In light-cone quantization with $A^+=0$, the $\overline{\ensuremath{\operatorname{MS}}}$ definition has an interpretation as a parton number density.

So far our main knowledge of the PDF is obtained from global fits to deep inelastic scattering and jet data, see for example~\cite{Ball:2014uwa,Dulat:2015mca,Martin:2009iq,Alekhin:2017kpj,Buckley:2014ana}. On the other hand, calculating the PDF from first principles with QCD has been an attractive subject, which can for example provide access to spin and momentum distributions that are hard to determine experimentally. Several different approaches to this have been considered using the lattice theory which is a nonperturbative method to solve QCD. Since the lattice theory is defined in a discretized Euclidean space with imaginary time, it is very difficult to calculate Minkowskian quantities with real-time dependence such as the PDF. The first and most well explored option is calculating the moments of the PDF~\cite{Martinelli:1987zd,Martinelli:1988xs,Detmold:2001dv,Detmold:2002nf,Dolgov:2002zm} that are matrix elements of local gauge-invariant operators. However, since the lattice regularization breaks $O(4)$ rotational symmetry, the consequent mixing between operators of different dimensions makes it difficult to compute higher moments, which in practice has limited the amount of information that can be extracted from this approach. A method to improve this situation by restoring the rotational symmetry has been proposed in Ref.~\cite{Davoudi:2012ya}.
Other proposals include extracting the PDF from the hadronic tensor~\cite{Liu:1993cv,Liu:1998um,Liu:1999ak,Liu:2016djw} and the forward Compton amplitude~\cite{Chambers:2017dov}, possibly with flavor changing currents~\cite{Detmold:2005gg}, and the more general ``lattice cross sections"~\cite{Ma:2014jla,Ma:2017pxb}. Systematic lattice analyses of these approaches are  under investigation, but challenges remain.

In Ref {\cite{Ji:2013dva}} Ji proposed that the $x$-dependence of the PDF can be extracted from a Euclidean distribution on the lattice, which can be understood in the language of the large momentum effective theory (LaMET)~\cite{Ji:2014gla}. This Euclidean distribution is referred to as the quasi-PDF, whose bare matrix element is defined using a spatial correlation of quarks along the $z$ direction,
\beq \label{eq:qpdf}
	\tilde{q}(x,P^z,\epsilon) \equiv \int_{-\infty}^\infty {dz\over 4\pi} e^{ixP^zz}\langle P | \bar{\psi}  (z) \Gamma U (z, 0) \psi (0) | P \rangle \,,
\eeq
where $\Gamma=\gamma^z$, $z^\mu=ze^\mu$, $e^\mu=(0,0,0,1)$, and the Wilson line is
\beq
	U(z,0) = P\exp\left(-ig\int_0^z dz' A^z(z')\right)\ .
\eeq
For finite momentum $P^z$, $\tilde{q}(x,P^z,\epsilon)$ has support in $-\infty < x < \infty$.
According to Ref.~\cite{Hatta:2013gta}, there is a universality class of operators that can be considered. For example, for the quasi-PDF, one could also replace $\Gamma=\gamma^z$ by $\Gamma=\gamma^0$ in Eq.~(\ref{eq:qpdf}) as both definitions reduce to the PDF under an infinite Lorentz boost along the $z$ direction. Unlike the PDF in Eq.~(\ref{eq:lcpdf}) that is invariant under the Lorentz boost, the quasi-PDF depends dynamically on it through the nucleon momentum $P^z$. When the nucleon momentum $P^z$ is much larger that the nucleon mass $M$ and $\Lambda_{\rm QCD}$, which is an attainable window on the lattice, the quasi-PDF can be factorized
into a matching coefficient and the PDF~\cite{Ji:2013dva,Ji:2014gla}. The factorization formula is
\begin{align} \label{eq:momfact}
	\tilde{q}(x,P^z,\mu_R)
	=& \int_{-1}^1 {dy\over |y|}\: C\Bigl({x\over y},{\mu_R\over \mu},{\mu\over p^z}\Bigr) \, q(y,\mu)
	\\
	& + \mathcal{O}\biggl({M^2\over P_z^2},{\Lambda_{\ensuremath{\operatorname{QCD}}}^2\over P_z^2}\biggr)\ ,\nn
\end{align}
where the renormalized quasi-PDF $\tilde{q}(x,P^z,\mu_R)$ is defined in a particular scheme at renormalization scale $\mu_R$, and $\mathcal{O}(M^2/P_z^2, \Lambda_{\rm QCD}^2/P_z^2)$ are power corrections suppressed by the nucleon momentum. In general the result for $C$ will depend on the choice of $\Gamma=\gamma^z$ or $\gamma^0$ and
renormalization schemes. For $\Gamma=\gamma^z$ the matching coefficient $C$ has been computed for the iso-vector quark quasi-PDF at one-loop level, first with a transverse momentum cutoff in~\cite{Xiong:2013bka}, confirmed in~\cite{Ma:2014jla,Alexandrou:2015rja}, and also recently determined in the regularization-invariant momentum subtraction (RI/MOM) scheme~\cite{Stewart:2017tvs}. Matching for the gluon quasi-PDF is calculated in~\cite{Wang:2017qyg,Wang:2017eel}. The matching coefficient $C$ is independent of the choice of states used for $\tilde q$ and $q$.\footnote{The scheme definition itself may separately involve a choice of state, such as in the RI/MOM scheme, but the result is still an operator renormalization that can be used for different choices of hadronic states for $\tilde q$ and $q$. The transverse cutoff and $\overline{\rm MS}$ schemes do not require even this choice.} Since matching calculations are carried out with quark states of momentum $p^z$, it can be tricky to know what the right choice to make is for $C$, and in some of the literature the choice of $p^z=P^z$ has been suggested when utilizing $C$ for the hadronic nucleon state. This is for example the case in the original quasi-PDF papers~\cite{Ji:2013dva,Ji:2014gla,Xiong:2013bka} and in the pioneering lattice calculations of the PDF from the quasi-PDF in~\cite{Lin:2014zya,Alexandrou:2015rja,Chen:2016utp,Alexandrou:2016jqi,Zhang:2017bzy,Alexandrou:2017huk,Chen:2017mzz,Lin:2017ani,Chen:2017gck}, which was summarized in Ref.~\cite{Lin:2017snn}. In the quasi-generalized parton distribution analysis in~\cite{Ji:2015qla} it was observed that one should take $p^z=|y|P^z$.  Through our rigorous analysis of \eq{momfact} we show that the correct result for this equation is indeed $p^z=|y|P^z$.

Recently, a different procedure~\cite{Radyushkin:2017cyf} to extract PDFs from the same lattice QCD matrix element as in~\cite{Ji:2013dva}
has been proposed based on the Lorentz invariant variables of the \spcorr (or pseudo-PDF), in place of the quasi-PDF. In this approach, one starts from the \spcorr $\tilde{Q}_{\gamma^\mu}$ defined for $\mu=0$ or $\mu=z$ by
\begin{align}
	{1\over2}\langle P | \tilde{O}_{\gamma^\mu}(z,\epsilon) | P\rangle
	= & P^\mu \tilde{Q}_{\gamma^\mu}(\zeta=P^zz,z^2,\epsilon) \,,
\end{align}
which depends on the two Lorentz invariants $z^2$ and $\zeta=-z\cdot P = P^z z$; the latter is also called the Ioffe time.
For an arbitrary Dirac matrix $\Gamma$ the operator $\tilde{O}_\Gamma$ is defined as
\begin{align} \label{eq:OGamma}
	\tilde{O}_{\Gamma}(z,\epsilon)  = \bar{\psi}  (z) \Gamma U (z, 0) \psi (0) \,.
\end{align}
This is the same spatial correlator (calculable on lattice) used to define the quasi-PDF in \eq{qpdf}, where $P^z$ is fixed and one Fourier transforms with respect to $z$. If instead $z^2$ is fixed, and we Fourier transform from the Ioffe time $\zeta$---which is in principle integrating over $P^z$---to the momentum fraction $x$, then one obtains the pseudo-PDF~\cite{Radyushkin:2017cyf},
\begin{align} \label{eq:pseudo}
\mathcal{P} \left( x, z^2,\epsilon\right)
   = \int_{-\infty}^\infty \! \frac{d \zeta}{2\pi}\:
    e^{ix \zeta}\: \tilde{Q}_{\gamma^0}\bigl( \zeta, z^2,\epsilon \bigr)
   \,.
\end{align}
For arbitrary finite $z$, the pseudo-PDF only has support in $-1\le x \le 1$~\cite{Radyushkin:1983wh,Radyushkin:2016hsy}, but has no parton model interpretation.
(Again the pseudo-PDF can equally well be considered for $\Gamma=\gamma^z$.) The \spcorr or pseudo-PDF approach has been explored on the lattice~\cite{Orginos:2017kos,Karpie:2017bzm}, where the short distance behavior was explored.
The PDF corresponds to the situation when $z^\mu$ is light-like, in which case the space-time correlator is referred to as the Ioffe-time distribution~\cite{Ioffe:1969kf},
\begin{align}
 Q(\zeta,\epsilon) = \tilde Q_{\gamma^+}(\zeta=-P^+\xi^-,z^2=0,\epsilon)
  \,.
\end{align}
When Fourier transformed this correlation gives the PDF
\begin{align} \label{eq:Qdefn}
q(y,\epsilon) =\int_{-\infty}^{\infty} {d\zeta\over 2\pi} \ e^{iy\zeta}\: Q(\zeta,\epsilon)\,.
\end{align}
In short the quasi-PDF and pseudo-PDF are different representations of the Euclidean \spcorr, as summarized in Table.~\ref{tab:I}.
\begin{table}
\begin{center}
	\begin{tabular}{ | c | c | c |}
		\hline
		Distribution & Fourier transform & Arguments\\[-2pt]
		& from \spcorr & \\ \hline
		Spatial correlation  &  & $\zeta = zP^z, z^2$\\ \hline
		Quasi-PDF & $z\to xP^z$ & $x, P^z$\\ \hline
		Pseudo-PDF & $\zeta \to x$ & $x, z^2$\\ \hline
	\end{tabular}
\caption{Summary of the relationship between different Euclidean distributions.}
\label{tab:I}
\end{center}
\end{table}

It was pointed out in Ref.~\cite{Ji:2017rah} that to obtain enough information to extract the PDF for the \spcorr with small $z^2$, one has to do lattice calculations with large momenta $P^z$, which is the same requirement as for the quasi-PDF. Ref.~\cite{Ji:2017rah} also proposed that the renormalized pseudo-PDF satisfies the following small $z^2$ factorization,
\begin{align} \label{eq:pseudofact}
	\mathcal{P} \left( x, z^2\mu_R^2\right)
  &= \int {dy\over |y|}\: \mathcal{C}\left({x\over y}, {\mu_R^2\over \mu^2},z^2\mu_R^2\right) q(y,\mu)
   \\
  &\ + {\cal O}(z^2\Lambda_{\rm QCD}^2, z^2 M^2) \,,\nn
\end{align}
which they verified at order $O(\alpha_s)$ for the unpolarized iso-vector case with $\Gamma=\gamma^0$.  (Again the coefficient ${\cal C}$ will depend on the choice of $\Gamma=\gamma^0$ or $\gamma^z$.)

In Ref.\cite{Ma:2014jla} a diagrammatic derivation of the factorization formula in \eq{momfact} for the quasi-PDF was given.  Here we derive this factorization formula for the quasi-PDF in an alternate manner, and also show that \spcorr and pseudo-PDF are different representations of the same
fundamental factorization. Our approach is based on the operator product expansion (OPE) for spacelike separated local operators~\cite{Wilson:1969zs}. For such operators the OPE has been proven for scalar field theory to all orders in perturbation theory~\cite{Brandt:1967rb,Wilson:1972ee,Zimmermann:1972tv}, and is widely assumed to hold for any renormalizable quantum field theory including QCD. By introducing auxiliary fields in place of the Wilson line~\cite{Dorn:1986dt}, the correlator in \eq{OGamma} is known to be equivalent to a product of local renormalizable operators of this type. Through our derivation we find the explicit form of the large $P^z$ and small $z^2$ factorization formulas in Eq.(\ref{eq:momfact}) and Eq.~(\ref{eq:pseudofact}) respectively, as well as the relationship between the matching coefficients $C$ and $\mathcal{C}$. Since the requirement for large $P^z$ and small $z^2$ is the same for both the quasi-PDF and pseudo-PDF approaches, there is in principle no fundamental difference in applying either one
to lattice calculations of the proton matrix element of $\tilde{O}_{\Gamma}(z)$. It is interesting to compare
both approaches utilizing the same lattice data, although they shall not yield different result in principle.

The rest of this paper is organized as follows: In Sec.~\ref{sec:ope}, we use an OPE of $\tilde{O}_\Gamma(z)$ to derive the large $P^z$ factorization of LaMET in Eq.~(\ref{eq:momfact}) and small $z^2$ factorization of the pseudo-PDF in Eq.~(\ref{eq:pseudofact}). We prove that one must take $p^z=|y|P^z$ in \eq{momfact}, so the corresponding argument in $C$ is $\mu/(|y|P^z)$.
(This OPE approach was used recently in Ref.~\cite{Ma:2017pxb} to prove the factorization theorem for the ``lattice cross sections", and the OPE proof carried out here was done independently and first presented in Ref.~\cite{Zhao:2017int}.)
%
In Sec.~\ref{sec:equiv}, we derive the \spcorr, pseudo-PDF, and quasi-PDF distributions and matching coefficients at one-loop in $\overline{\rm MS}$ and analyze the Fourier-transform relation between the quasi-PDF and pseudo-PDF. Unlike earlier results for the quasi-PDF in $\overline{\rm MS}$, we also use dimensional regularization with minimal subtraction to renormalize divergences at $x=\pm\infty$. In Sec.~\ref{sec:ren}, we discuss how renormalization schemes other than $\overline{\rm MS}$ are easily incorporated into the factorization formulas.
%
In Sec.~\ref{sec:num} we carry out a numerical analysis of the matching coefficients, by computing the convolution in Eq.~(\ref{eq:momfact}) numerically using the PDF determined by global fits~\cite{Martin:2009iq}. We show that the difference between using $p^z=P^z$ and $p^z=|y|P^z$ in \eq{momfact} is an important effect, and that our $\overline{\rm MS}$ matching coefficients are insensitive to cutoffs in the convolution integral. In Sec.~\ref{sec:lattice}, we discuss the implications of our OPE analysis for the lattice calculation of the PDF in both the quasi-PDF and pseudo-PDF approaches. Finally, we conclude in Sec.~\ref{sec:summary}.
